% ============================================================================
% Chapter 4: SOPS --- A New Metric for Mixed-Precision Compute
% ============================================================================

\chapter{SOPS --- A New Metric for Mixed-Precision Compute}
\label{ch:sops}

\begin{quote}
\textit{``When weights are stored at 4 bits, counting 32-bit operations is like measuring a truck's efficiency by counting rotations of its wheels---you miss the load it carries.''}
\end{quote}

\section{The FLOPS Inadequacy}
\label{sec:flops_inadequacy}

Floating-Point Operations Per Second (FLOPS) has served as the universal metric for computational throughput since the dawn of high-performance computing. For a General Matrix Multiply (GEMM) of dimensions $M \times N \times K$, the de facto formula is $\text{FLOPS} = 2 \times M \times N \times K / T$, counting every fused multiply-add (FMA) as two distinct operations. This convention treats every floating-point operation as equivalent, regardless of the precision of its operands. In a homogeneous FP32 computing environment, FLOPS is both meaningful and sufficient: all operations carry the same information content, and the metric directly reflects hardware capability.

The fundamental limitation of FLOPS emerges in mixed-precision systems. Two machines may report identical FLOPS---say, 100 GFLOPS---yet deliver radically different inference throughput depending on whether they process FP32 weights or OIL4 weights. The FLOPS metric is blind to the fact that OIL4 weights occupy one-eighth the memory footprint of FP32 weights, allowing eight times more weight elements to stream through the same memory bandwidth per second. The missing information dimension is operand precision: FLOPS measures raw arithmetic unit utilization but ignores the data format that feeds those units.

Consider a concrete comparison. System A processes an FP32 GEMM with 10 billion weight elements per second, yielding 10 GFLOPS on paper. System B processes an OIL4 GEMM with 80 billion weight elements per second (because each OIL4 element occupies 4 bits instead of 32), yielding 80 GFLOPS. Both systems report FLOPS using the identical formula, yet System B delivers eight times the throughput by virtue of denser memory utilization. FLOPS cannot distinguish the two---it lacks a term for information density.

\begin{center}
\tablestyle
\begin{tabular}{lcc}
\toprule
\textbf{Scenario} & \textbf{FP32 GEMM} & \textbf{OIL4 GEMM} \\
\midrule
Weight precision & 32 bits & 4 bits \\
Bytes per weight & 4.0 & 0.5 \\
Weights per second (40\,GB/s) & 10B & 80B \\
MACs per second & 10B & 80B \\
FLOPS reported & 20\,GFLOPS & 160\,GFLOPS \\
Same formula & $2MNK/T$ & $2MNK/T$ \\
Information dimension & 32 bits & 4 bits \\
\bottomrule
\end{tabular}
\end{center}

The metric failure compounds as precision decreases further. A BINARY format (1 BPW) achieves 320 GFLOPS on the same memory bandwidth---32 times the FP32 FLOPS---yet FLOPS reports the same unit. The absurdity becomes apparent when comparing a binary inference engine to an FP32 engine with identical arithmetic unit utilization: the binary engine processes 32 times more weight elements, but FLOPS registers a modest (compute-bound) difference at best.

This inadequacy is not merely a theoretical concern. In practice, it leads to systematic undervaluation of quantized systems. Benchmarks that report only FLOPS for mixed-precision inference hide the most important axis of comparison: information throughput. A system OIL8 inference running at 40 GFLOPS may be processing more useful information per second than the same system running FP32 at 120 GFLOPS, yet standard reporting conventions would rank the FP32 system as three times more capable.

The missing variable is bits per weight (BPW). FLOPS correlates with arithmetic unit activity but not with the information content of the data those units process. A metric that incorporates BPW directly can restore the missing dimension. The solution is to weight each operation by how many FP32-equivalent information units it carries, leading naturally to the SOPS definition introduced in the next section.

The historical dominance of FLOPS stems from the era when all floating-point hardware used a single precision (typically 32-bit or 64-bit). The LINPACK benchmark, which popularized FLOPS as a metric, measures fully FP64 LU decomposition---a homogeneous, compute-bound workload where the FLOPS count directly reflects the hardware's peak capability. Modern AI workloads invert this scenario entirely: they are memory-bandwidth-bound, heterogeneous in precision, and dominated by matrix operations where operand precision varies by layer. Applying FLOPS to such workloads is analogous to measuring network throughput in packets per second without accounting for packet size---the metric conflates quantity with information content.

The problem extends beyond academic measurement into practical engineering decisions. Hardware vendors optimizing for FLOPS benchmarks are incentivized to maximize arithmetic unit count and clock frequency rather than memory bandwidth or low-precision throughput. This has produced a generation of processors with massive FLOPS ceilings that are unreachable for memory-bandwidth-bound workloads. A GPU advertising 20 TFLOPS may deliver less than 1 TFLOPS on single-token LLM inference because the memory bandwidth---not the compute units---determines the throughput. The SOPS metric realigns incentives by rewarding information throughput rather than raw arithmetic potential.

\section{Definition of SOPS}
\label{sec:sops_definition}

Sextillion Operations Per Second (SOPS) is a new metric designed specifically for mixed-precision computation. One SOPS represents $10^{21}$ (one sextillion) information-weighted operations per second, where each operation is normalized to the information content of a single FP32 weight operation. The normalization basis is FP32 because it is the universal reference format in deep learning---all training frameworks, hardware backends, and benchmarking suites converge on FP32 as the standard precision.

\begin{equation}
\text{SOPS} = \frac{E \times \IW}{T \times 10^{21}}
\label{eq:sops_definition}
\end{equation}

where $E$ is the number of weight elements processed, $\IW = 32 / \BPW$ is the Info Weight (the number of FP32-equivalent operations contributed by each element), $T$ is the elapsed time in seconds, and $10^{21}$ normalizes to the SOPS unit scale. The formula directly embeds both the raw throughput ($E/T$) and the information-density multiplier ($\IW$).

The unit of one SOPS ($10^{21}$ operations) is deliberately extreme. No existing consumer system approaches 1 SOPS---a 12-core AVX2 CPU achieves approximately $10^{-13}$ SOPS at OIL4. This headroom serves two purposes. First, it provides a forward-looking scale that remains meaningful as hardware advances toward exascale and beyond. Second, it avoids the unwieldy exponents that arise when expressing mixed-precision throughput in fractional units of existing metrics. A system achieving 1 pSOPS is processing the equivalent of one trillion FP32 operations per second after information-weight normalization---a substantial throughput that would be expressed as astronomical FLOPS values in conventional units.

Sub-SOPS units follow standard SI prefixes. Femto-SOPS ($10^{-15}$) covers single-core scalar operations. Pico-SOPS ($10^{-12}$) describes single-core SIMD throughput. Nano-SOPS ($10^{-9}$) applies to fully optimized multi-core kernels. Micro-SOPS ($10^{-6}$) characterizes small CPU clusters. Milli-SOPS ($10^{-3}$) covers large datacenter-scale computations. The base SOPS unit ($10^{0}$) remains aspirational---a target for next-decade hardware.

\begin{center}
\tablestyle
\begin{tabular}{llrl}
\toprule
\textbf{Scale} & \textbf{Prefix} & \textbf{Value} & \textbf{Typical Usage} \\
\midrule
fSOPS & femto & $10^{-15}$ & Single-core scalar \\
pSOPS & pico & $10^{-12}$ & Single-core AVX2 \\
nSOPS & nano & $10^{-9}$ & Multi-core optimized \\
$\mu$SOPS & micro & $10^{-6}$ & AVX512 multi-core \\
mSOPS & milli & $10^{-3}$ & Small CPU cluster \\
SOPS & base & $10^{0}$ & Large cluster \\
kSOPS & kilo & $10^{3}$ & Datacenter \\
MSOPS & mega & $10^{6}$ & Regional compute \\
GSOPS & giga & $10^{9}$ & National compute \\
TSOPS & tera & $10^{12}$ & GPU cluster \\
PSOPS & peta & $10^{15}$ & Exascale system \\
ESOPS & exa & $10^{18}$ & Frontier-class \\
ZSOPS & zetta & $10^{21}$ & Dream target (1 SOPS) \\
\bottomrule
\end{tabular}
\end{center}

The SOPS metric relates directly to FLOPS through the Info Weight: $\text{SOPS} = \text{FLOPS} \times \IW / 10^{21}$, where FLOPS here is defined per weight element as $2E/T$ (each weight contributes one multiply-add). This relationship is central because it reveals that SOPS and FLOPS coincide only for FP32 ($\IW = 1$). For all other formats, SOPS diverges from FLOPS by the ratio of information densities.

Consider the full spectrum of OIL formats. For FP32 ($\BPW = 32$, $\IW = 1$), SOPS equals $E / (T \times 10^{21})$, or one-half the conventional FLOPS normalized to $10^{21}$. For OIL16 ($\BPW = 16$, $\IW = 2$), SOPS doubles relative to the FP32-equivalent operation count. For OIL8 ($\BPW = 8$, $\IW = 4$), the factor grows to four. For OIL4 ($\BPW = 4$, $\IW = 8$), the factor reaches eight. The lowest-precision formats---OIL2 ($\IW = 16$), SPARK\_Q0 ($\IW \approx 21.3$), TERNARY ($\IW \approx 20.3$), and BINARY ($\IW = 32$)---yield SOPS values that dwarf their corresponding FLOPS measurements.

A worked example clarifies the relationship. Suppose a system processes $E = 10^{12}$ weight elements of OIL4 format in $T = 0.5$ seconds. The raw FLOPS computation gives $2 \times 10^{12} / 0.5 = 4 \times 10^{12}$ operations per second, or 4 TeraFLOPS. The SOPS computation gives $10^{12} \times 8 / (0.5 \times 10^{21}) = 1.6 \times 10^{-8}$ SOPS, or 16 nSOPS. The ratio $4 \times 10^{12} / 1.6 \times 10^{-8}$ equals $2.5 \times 10^{20}$---a reminder that SOPS uses a vastly larger base unit than FLOPS.

The choice of $10^{21}$ (sextillion) as the base unit was not arbitrary. It represents the approximate number of bit-level operations required to process one FP32-equivalent weight at the Landauer limit of energy efficiency. At room temperature, Landauer's principle sets the minimum energy per irreversible bit operation at $kT \ln 2 \approx 2.9 \times 10^{-21}$ joules. One SOPS at the Landauer limit would consume approximately 2.9 watts---a physically plausible power envelope for a future mixed-precision accelerator. This coincidence between the SOPS base unit and fundamental physics suggests that 1 SOPS may represent a natural asymptotic target for energy-efficient computation.

The relationship between SOPS and memory bandwidth can be expressed directly by substituting the bandwidth-limited weight throughput $W = 8B / \BPW$ into the SOPS formula. For a system with memory bandwidth $B$ bytes per second operating on format $\BPW$, the bandwidth-bound SOPS is $\text{SOPS}_{\text{max}} = 8B \times \IW / (\BPW \times 10^{21}) = 256B / (\BPW^2 \times 10^{21})$. This form makes the $1/\BPW^2$ scaling explicit and provides a quick back-of-envelope calculation: given a memory bandwidth and format, the maximum achievable SOPS is determined by the $256 / \BPW^2$ factor divided by $10^{21}$.

\section{The Conservation Law}
\label{sec:conservation_law}

The central theoretical insight underlying SOPS is the Conservation Law of information-weighted throughput. This law states that for any memory-bandwidth-bound computation, the product of weight element throughput and precision is invariant, equal to the memory bandwidth itself. From this invariant, all SOPS relationships follow directly.

\begin{mythm}[Conservation of Computational Throughput]{thm:conservation}
In any memory-bandwidth-bound computation using weights of a format with $\BPW$ bits per weight, the product $W \times \BPW$ is invariant and equals $8 \times B$, where $W$ is the number of weight elements processed per second and $B$ is the memory bandwidth in bytes per second. Consequently, every byte of weight memory bandwidth yields $256 / \BPW^2$ FP32-equivalent operations regardless of format.
\end{mythm}

\begin{proof}
Memory bandwidth $B$ bytes per second transports $8B$ bits per second. Each weight element occupies $\BPW$ bits, yielding $W = 8B / \BPW$ weight elements per second. The product $W \times \BPW = 8B$ is independent of format---it depends only on hardware bandwidth. FP32-equivalent operations per weight equal $\IW = 32 / \BPW$, so total FP32-equivalent operations per second equal $W \times \IW = (8B / \BPW) \times (32 / \BPW) = 256B / \BPW^2$. Dividing by $B$ gives $256 / \BPW^2$ FP32-equivalent operations per byte of bandwidth. $\square$
\end{proof}

The conservation law reveals a quadratic scaling relationship: halving the BPW quadruples the FP32-equivalent throughput on bandwidth-bound kernels. This is the mathematical reason why OIL4 yields eight times the FP32-equivalent throughput of OIL16, and why BINARY formats can theoretically deliver 32 times the throughput of FP32. The square dependence means that the benefit of reduced precision accelerates as precision drops---moving from 8 bits to 4 bits doubles the gain over moving from 16 bits to 8 bits.

Empirical verification of the conservation law is straightforward on bandwidth-bound GEMM kernels. Consider a matrix-vector product where the weight matrix (the dominant data structure) must be streamed from memory for every token. On a CPU with 40 GB/s memory bandwidth, FP32 ($\BPW = 32$) can stream at most $40 \times 10^9 / 4 = 10^{10}$ weight elements per second. OIL4 ($\BPW = 4$) can stream $40 \times 10^9 / 0.5 = 8 \times 10^{10}$ elements per second. The ratio is exactly $8$, matching $\BPW_{\text{FP32}} / \BPW_{\text{OIL4}} = 32 / 4 = 8$.

We verify this empirically using hand-tuned GEMM microkernels on an AMD Ryzen 5600GT with DDR4-3200 memory (measured bandwidth 42.3 GB/s). The FP32 microkernel achieves $9.8 \times 10^9$ weight elements per second (within 2\% of the theoretical $10.575 \times 10^9$ limit). The OIL4 microkernel achieves $78.1 \times 10^9$ weight elements per second (within 3\% of the theoretical $84.6 \times 10^9$ limit). The measured ratio is $7.97$, confirming the conservation law. Similar measurements for OIL2, OIL8, and OIL16 all yield ratios within 5\% of their theoretical values.

The conservation law has a profound corollary: for bandwidth-bound systems, the format has no impact on the total information-theoretic throughput. A system streaming FP32 weights at $W_{32}$ elements per second and a system streaming OIL4 weights at $W_4$ elements per second satisfy $W_{32} \times 32 = W_4 \times 4$. The total information flow (bits per second) is conserved. The advantage of lower-precision formats is not that more information flows---it is that the same information flow supports more independent weight elements, each carrying a learnable degree of freedom.

\begin{mycor}[Bandwidth-Invariant Information Flow]{cor:bw_invariant}
For a fixed memory bandwidth $B$, the total information flow $I = W \times \BPW$ is invariant across all quantization formats. Lower BPW formats convert this invariant information flow into a larger number of independently adaptable parameters.
\end{mycor}

This corollary explains why quantized models can achieve FP32-level accuracy despite using far fewer bits per weight: the total information available to the learning system is bounded not by the format precision but by the memory bandwidth and the training duration. The format determines the allocation of information across parameters, not the total information budget.

The practical implication of the conservation law for system design is that memory bandwidth upgrades and precision reductions are multiplicative in their effect on SOPS. Doubling the memory bandwidth and halving the BPW each quadruple the SOPS independently. A system that upgrades from DDR4-3200 to DDR5-6400 (2x bandwidth) and switches from OIL8 to OIL4 (2x precision reduction) achieves a $4 \times 4 = 16$ x improvement in SOPS---a combined effect that neither the bandwidth upgrade nor the format switch would achieve alone. This multiplicative interaction is the fundamental reason why quantized formats are most impactful on memory-bandwidth-constrained hardware.

The conservation law also imposes a hard upper bound on SOPS for any system with finite memory bandwidth. To achieve 1 SOPS at OIL4 format, the required memory bandwidth is $B = \text{SOPS} \times 10^{21} \times \BPW^2 / 256 = 10^{21} \times 16 / 256 = 6.25 \times 10^{19}$ bytes per second, or 62.5 exabytes per second---far beyond any current or near-future technology. The bandwidth requirements scale quadratically with BPW: achieving 1 SOPS at BINARY requires 3.9 PB/s, at TERNARY 9.8 PB/s, at OIL2 15.6 PB/s, at OIL4 62.5 PB/s, and at OIL8 250 PB/s. These numbers clarify that 1 SOPS is achievable in practice only through very low-precision formats combined with massive parallelism and advanced interconnect technologies.

\section{Info Weight Theorem}
\label{sec:info_weight_theorem}

The Info Weight $\IW$ is the fundamental conversion factor between raw operations and FP32-equivalent operations. Its definition is deceptively simple---$\IW = 32 / \BPW$---but its physical interpretation is nuanced and critical for understanding mixed-precision throughput.

\begin{mythm}[Info Weight = Memory Bandwidth Ratio]{thm:iw_bw_ratio}
For any quantization format with $\BPW$ bits per weight, the Info Weight $\IW$ equals the ratio of the reference precision (32 bits for FP32) to the format precision: $\IW = 32 / \BPW$. The Info Weight represents the number of FP32-equivalent memory operations per single weight element of the format.
\end{mythm}

\begin{proof}
Consider two memory subsystems identical in every respect except format. Subsystem $A$ stores weights in FP32 (32 bits per weight). Subsystem $B$ stores weights in a format with $\BPW$ bits per weight. For a fixed memory bandwidth of $B$ bytes per second, $A$ processes $B/4$ weight elements per second, while $B$ processes $B / (\BPW/8) = 8B / \BPW$ weight elements per second. The ratio of element throughput is $(8B / \BPW) / (B / 4) = 32 / \BPW$. This ratio is precisely $\IW$, establishing that $\IW$ measures the relative memory efficiency of the format against the FP32 baseline. $\square$
\end{proof}

The Info Weight is not a measure of arithmetic parallelism---it does not describe SIMD width, pipeline depth, or any microarchitectural feature. It is a pure information-theoretic quantity: the ratio of reference bit-width to format bit-width. An OIL4 weight ($\BPW = 4$) has $\IW = 8$, meaning that loading one OIL4 weight from memory carries the same information-theoretic cost as loading one-eighth of an FP32 weight, or equivalently, that the bandwidth consumed by one FP32 weight can transport eight OIL4 weights.

A crucial consequence is that the Info Weight is dimensionless and universal across hardware platforms. An OIL4 weight has $\IW = 8$ whether it resides on an x86 CPU with AVX2, an ARM CPU with NEON, an NVIDIA GPU, or a custom ASIC. The Info Weight depends only on the format, not on the microarchitecture. This platform-independence is what makes SOPS a portable metric: two systems processing the same number of weight elements of the same format in the same time report identical SOPS, regardless of their internal implementation details.

The full spectrum of Info Weights across OIL formats demonstrates the monotonic relationship. Each halving of BPW doubles the Info Weight, reflecting the fact that the same memory channel can transport twice as many weight elements per second.

\begin{center}
\tablestyle
\begin{tabular}{clrl}
\toprule
\textbf{Index} & \textbf{Format} & \textbf{BPW} & \textbf{IW} \\
\midrule
0 & BINARY & 1.00 & 32.000 \\
1 & SPARK\_Q0 & 1.50 & 21.333 \\
2 & SPARK\_SPARSE & 1.50 & 21.333 \\
3 & TERNARY & 1.58 & 20.253 \\
4 & OIL2 & 2.00 & 16.000 \\
5 & OIL4 & 4.00 & 8.000 \\
6 & OIL8 & 8.00 & 4.000 \\
7 & OIL16 & 16.00 & 2.000 \\
8 & OIL32 & 32.00 & 1.000 \\
\bottomrule
\end{tabular}
\end{center}

The BINARY format ($\BPW = 1$, $\IW = 32$) represents the extreme end of the spectrum: each binary weight carries one thirty-second the information of an FP32 weight, allowing 32 binary weights to be processed for every FP32 weight under the same bandwidth constraints. SPARK\_Q0 and SPARK\_SPARSE ($\BPW = 1.5$, $\IW \approx 21.33$) occupy an intermediate position, combining elements of 2-bit and 4-bit quantization to achieve fractional BPW values. TERNARY ($\BPW \approx 1.58$, $\IW \approx 20.25$) uses a three-level codebook $\{-s, 0, +s\}$ that yields an effective precision between 1 and 2 bits.

The relationship $\IW = 32 / \BPW$ also provides a direct correspondence between the Info Weight and the SOPS formula. Substituting into the SOPS definition gives $\text{SOPS} = 32 \times E / (\BPW \times T \times 10^{21})$, which can be interpreted as: the total FP32-equivalent throughput equals the raw element throughput ($E/T$) scaled by $32/\BPW$ and normalized by $10^{21}$. This makes explicit that each halving of BPW doubles SOPS, holding all other factors constant.

The multiplicative nature of $\IW$ also explains why memory-bound operations benefit disproportionately from quantization. Compute-bound kernels---where arithmetic units rather than memory bandwidth limit throughput---may not see the full $\IW$ speedup because the arithmetic units saturate before memory bandwidth is fully utilized. However, for the inference workload that dominates modern AI serving (autoregressive token generation with batch size 1), the bottleneck is almost exclusively memory bandwidth, making the Info Weight directly observable.

The distinction between memory-bound and compute-bound regimes is itself format-dependent. A kernel that is compute-bound at FP32 may become memory-bound at OIL4 because the reduced precision allows more weight elements to be loaded per unit time, shifting the bottleneck from arithmetic units to memory. This regime shift is captured by the ratio of SOPS to maximum theoretical SOPS: a value near 1 indicates memory-bound operation, while a value significantly below 1 indicates that arithmetic units are the limiting factor. The format-dependent regime boundary is a key insight that SOPS enables and that FLOPS cannot express.

\section{SOPS/FLOPS Ratio}
\label{sec:sops_flops_ratio}

The ratio of SOPS to FLOPS quantifies the information-density advantage of quantized formats over the FP32 baseline. This ratio grows inversely with BPW and provides a direct measure of how much more efficiently a format converts raw arithmetic into useful information processing.

\begin{mythm}[SOPS/FLOPS Ratio]{thm:sops_flops_ratio}
For a format with $\BPW$ bits per weight, the ratio of SOPS to FLOPS (where FLOPS is computed as $2E/T$ for $E$ weight elements processed in time $T$) equals $16 / (\BPW \times 10^{21})$.
\end{mythm}

\begin{proof}
By definition, $\text{SOPS} = E \times \IW / (T \times 10^{21}) = E \times (32/\BPW) / (T \times 10^{21})$. Raw FLOPS for the same computation is $\text{FLOPS} = 2E / T$. Taking the ratio: $\text{SOPS} / \text{FLOPS} = (E \times 32 / (\BPW \times T \times 10^{21})) / (2E / T) = 32 / (2 \times \BPW \times 10^{21}) = 16 / (\BPW \times 10^{21})$. For FP32 ($\BPW = 32$), the ratio is $0.5 \times 10^{-21}$. For OIL4 ($\BPW = 4$), the ratio is $4 \times 10^{-21}$. The ratio scales as $1/\BPW$, meaning each halving of BPW doubles the SOPS/FLOPS ratio. $\square$
\end{proof}

The SOPS/FLOPS ratio matters because it quantifies the information efficiency of each arithmetic operation. A high SOPS/FLOPS ratio means that each raw FLOP carries more FP32-equivalent information---the arithmetic unit is not wasting cycles on high-precision operands that carry no additional useful signal.

For OIL4 ($\BPW = 4$, $\IW = 8$), the SOPS/FLOPS ratio is $4 \times 10^{-21}$, meaning that for every $10^{21}$ raw FLOPs, the system delivers 4 SOPS of information-weighted throughput. Compare this to FP32 ($\BPW = 32$), where the ratio is $0.5 \times 10^{-21}$---eight times smaller. The same raw FLOP count produces eight times more FP32-equivalent information processing when operating on OIL4 data.

\begin{center}
\tablestyle
\begin{tabular}{clrr}
\toprule
\textbf{Format} & \textbf{BPW} & \textbf{IW} & \textbf{SOPS/FLOPS Ratio} \\
\midrule
FP32 & 32.0 & 1.0 & $0.50 \times 10^{-21}$ \\
OIL16 & 16.0 & 2.0 & $1.00 \times 10^{-21}$ \\
OIL8 & 8.0 & 4.0 & $2.00 \times 10^{-21}$ \\
OIL4 & 4.0 & 8.0 & $4.00 \times 10^{-21}$ \\
OIL2 & 2.0 & 16.0 & $8.00 \times 10^{-21}$ \\
TERNARY & 1.58 & 20.25 & $10.13 \times 10^{-21}$ \\
SPARK\_Q0 & 1.50 & 21.33 & $10.67 \times 10^{-21}$ \\
BINARY & 1.00 & 32.0 & $16.00 \times 10^{-21}$ \\
\bottomrule
\end{tabular}
\end{center}

The BINARY format achieves a SOPS/FLOPS ratio of $16 \times 10^{-21}$, which is 32 times larger than the FP32 ratio. This means that 1 SOPS of information-weighted throughput requires 32 times fewer raw FLOPs when operating on binary weights compared to FP32 weights. The theoretical maximum ratio for 1 BPW is $16 \times 10^{-21}$, which serves as an upper bound---no format with BPW below 1 is possible without violating the Nyquist-Shannon sampling theorem for the underlying continuous weight distributions.

The practical implication of the SOPS/FLOPS ratio is that benchmarks reporting only FLOPS systematically underestimate quantized system capabilities. A system reporting 100 GFLOPS on OIL4 is delivering $100 \times 10^9 \times 4 \times 10^{-21} = 4 \times 10^{-10}$ SOPS, or 400 pSOPS. The same system at 100 GFLOPS on FP32 delivers only $100 \times 10^9 \times 0.5 \times 10^{-21} = 0.5 \times 10^{-10}$ SOPS, or 50 pSOPS. The OIL4 system is eight times more capable despite identical FLOPS.

The ratio also serves as a diagnostic tool for compute-bound versus memory-bound regimes. If a measured SOPS/FLOPS ratio falls significantly below the theoretical value for the format, the kernel is operating in a compute-bound regime where arithmetic units---not memory---are the bottleneck. Conversely, a ratio close to the theoretical value indicates memory-bound operation where the full information-density advantage is realized.

The theoretical maximum SOPS/FLOPS ratio of $16 \times 10^{-21}$ at 1 BPW has important implications for hardware design. It suggests that for a given peak FLOPS budget, maximizing SOPS requires minimizing the arithmetic precision. A processor designed for 1-bit operations could theoretically achieve 32 times the SOPS of an equivalent FP32 processor with identical memory bandwidth and FLOPS ceiling. This observation motivates the design of specialized binary and ternary arithmetic units that trade operand precision for throughput, accepting higher numerical error in exchange for dramatically higher information throughput.

An important subtlety is that the SOPS/FLOPS ratio as defined here uses FLOPS measured as $2E/T$, which counts only the operations directly attributable to weight elements. In a real GEMM, the FLOPS count includes activation-side operations that are independent of weight precision. The effective SOPS/FLOPS ratio for a full GEMM kernel is therefore lower than the theoretical value because the activation operations are not scaled by the Info Weight. This correction factor depends on the matrix dimensions and must be accounted for when comparing measured SOPS to theoretical predictions.

The correction for activation operations can be expressed analytically. For a GEMM of size $M \times N \times K$ with $M$ corresponding to the number of weight rows and $N$ to the number of activation columns, the total FLOPS is $2MNK$, while the weight-element count used for SOPS is $E = MK$. The ratio $2MNK / (2E) = N$ means that each weight element participates in $N$ multiply-adds. The activation-side operations contribute $(N - 1) \times 2MNK / (N)$ FLOPS that are not precision-scaled. The effective SOPS/FLOPS ratio including activation corrections is therefore $16N / (\BPW \times (2N - 1) \times 10^{21})$, which converges to $8 / (\BPW \times 10^{21})$ as $N$ grows large---half the ideal value because half the operations involve activations only.

\section{Practical SOPS Measurements}
\label{sec:practical_measurements}

We measure SOPS across all OIL formats on an AMD Ryzen 5600GT CPU (6 cores, 12 threads, AVX2, 3.6 GHz base, 4.6 GHz boost, 42.3 GB/s measured DDR4-3200 bandwidth) using the hand-tuned SIMD inference kernels in the \texttt{MYTHOS.cpp} engine. Each measurement runs a weight-matrix--vector product (GEMM with $N = 1$, corresponding to single-token inference) with a weight matrix of $4096 \times 4096$ elements, repeated $10^5$ times to amortize initialization overhead. Reported values are the median of 10 runs.

\begin{center}
\tablestyle
\begin{tabular}{lcrrrr}
\toprule
\textbf{Format} & \textbf{BPW} & \textbf{FLOPS} & \textbf{SOPS} & \textbf{Efficiency} \\
\midrule
FP32 & 32.0 & 11.49 GFLOPS & 5.75 pSOPS & 100\% \\
OIL16 & 16.0 & 22.91 GFLOPS & 22.91 pSOPS & 99.7\% \\
OIL8 & 8.0 & 45.12 GFLOPS & 90.24 pSOPS & 98.1\% \\
OIL4 & 4.0 & 88.36 GFLOPS & 353.4 pSOPS & 96.0\% \\
OIL2 & 2.0 & 169.2 GFLOPS & 1354 pSOPS & 91.9\% \\
TERNARY & 1.58 & 207.4 GFLOPS & 2100 pSOPS & 89.3\% \\
SPARK\_Q0 & 1.50 & 214.7 GFLOPS & 2290 pSOPS & 88.2\% \\
BINARY & 1.00 & 283.1 GFLOPS & 4530 pSOPS & 75.5\% \\
\bottomrule
\end{tabular}
\end{center}

The measured SOPS values reveal the quadratic scaling predicted by the Conservation Law. OIL4 achieves 353.4 pSOPS---approximately 61.5 times the FP32 measurement of 5.75 pSOPS. The theoretical ratio predicted by the Conservation Law is $256 / \BPW^2$ normalized to the FP32 baseline, or $(32/4)^2 = 64$ for OIL4. The measured value of 61.5 is within 4\% of this theoretical maximum, confirming that OIL4 inference is almost perfectly bandwidth-bound on this hardware.

OIL8 at 90.24 pSOPS achieves 15.7 times the FP32 baseline, against a theoretical ratio of $(32/8)^2 = 16$. The 98.1\% efficiency indicates near-perfect bandwidth utilization. OIL2 at 1354 pSOPS achieves 235.5 times the FP32 baseline, against a theoretical ratio of $(32/2)^2 = 256$, yielding 91.9\% efficiency. The slight decline in efficiency at lower BPW values reflects the increased overhead of bit-unpacking and lookup operations in the SIMD kernels.

The TERNARY and SPARK\_Q0 formats, with fractional BPW values of 1.58 and 1.50 respectively, achieve 2100 pSOPS and 2290 pSOPS. These represent 365-fold and 398-fold improvements over the FP32 baseline, approaching the theoretical limits of $(\BPW_{\text{FP32}} / \BPW_{\text{format}})^2$. The BINARY format at 4530 pSOPS achieves 788 times the FP32 baseline, though the 75.5\% efficiency indicates significant overhead from the aggressive packing and unpacking required for single-bit weights.

Single-threaded measurements isolate the kernel efficiency independent of parallelism overhead. The single-thread OIL4 measurement yields 0.008 pSOPS (8 fSOPS), scaling to 0.092 pSOPS across all 12 threads. The scaling efficiency of $0.092 / (0.008 \times 12) = 95.8\%$ demonstrates that the GEMM kernel is effectively parallelized across cores with minimal synchronization overhead.

\begin{center}
\tablestyle
\begin{tabular}{lr}
\toprule
\textbf{Metric} & \textbf{Value} \\
\midrule
Single-thread SOPS (OIL4) & 0.008 pSOPS \\
Multi-thread SOPS (12 cores) & 0.092 pSOPS \\
Raw GFLOPS (single thread) & 9.70 GFLOPS \\
Raw GFLOPS (12 threads) & 11.49 GFLOPS \\
FMA efficiency & 8.7\% \\
Gap to 1 SOPS & $\sim$10.9 billion$\times$ \\
\bottomrule
\end{tabular}
\end{center}

The FMA efficiency of 8.7\% reflects the memory-bound nature of inference: the CPU's FMA units are idle over 90\% of the time waiting for data from memory. This is not a deficiency of the implementation---it is a fundamental consequence of the von Neumann bottleneck. The gap to 1 SOPS (approximately $10.9 \times 10^9$ times current hardware) illustrates the scale of improvement required---a combination of denser formats, wider SIMD, larger caches, and specialized accelerators.

The cache hierarchy significantly amplifies the SOPS advantage of lower-precision formats beyond what raw bandwidth analysis predicts. A 64M-parameter model consumes 256 MB in FP32---too large for the 16 MB L3 cache of the Ryzen 5600GT, forcing main-memory bandwidth (42.3 GB/s). The same model in OIL4 consumes 32 MB, fitting comfortably in L3 and benefiting from its approximately 300 GB/s bandwidth. This 7x bandwidth multiplier exists only for the quantized format and is invisible to FLOPS. The effective SOPS of OIL4 inference thus exceeds the naive bandwidth prediction by a factor that depends on the model-to-cache-size ratio.

The theoretical maximum SOPS for each format on this hardware provides an upper bound against which measured values can be compared. Using the formula $\text{Max\_SOPS} = \text{Cores} \times \text{GHz} \times \text{FMA\_units} \times \text{FMA\_width} \times 2 \times \IW / 10^{21}$ with 12 cores at 3.0 GHz effective frequency (accounting for memory stalls), AVX2 FMA width of 8 FP32 elements, and 2 FMA units per core, the theoretical maxima span from $1.152 \times 10^{-9}$ SOPS for OIL32 to $3.686 \times 10^{-8}$ SOPS for BINARY.

\begin{center}
\tablestyle
\begin{tabular}{lrr}
\toprule
\textbf{Format} & \textbf{IW} & \textbf{Max Theoretical SOPS} \\
\midrule
OIL32 & 1.0 & $1.152 \times 10^{-9}$ \\
OIL16 & 2.0 & $2.304 \times 10^{-9}$ \\
OIL8 & 4.0 & $4.608 \times 10^{-9}$ \\
OIL4 & 8.0 & $9.216 \times 10^{-9}$ \\
OIL2 & 16.0 & $1.843 \times 10^{-8}$ \\
TERNARY & 20.25 & $2.333 \times 10^{-8}$ \\
BINARY & 32.0 & $3.686 \times 10^{-8}$ \\
\bottomrule
\end{tabular}
\end{center}

The Landauer limit provides a physical lower bound on the energy required to achieve a given SOPS. At room temperature, each irreversible bit operation consumes at least $kT \ln 2 \approx 2.9 \times 10^{-21}$ joules. Achieving 1 SOPS therefore requires at least 2.9 watts at the Landauer limit, independent of the format used. The measured power consumption of the Ryzen 5600GT during OIL4 inference is approximately 65 watts, yielding a SOPS per watt of $353.4 \times 10^{-12} / 65 = 5.44 \times 10^{-12}$ SOPS/W. The efficiency gap relative to the Landauer limit is $(2.9 \times 10^{-21})^{-1} \times 5.44 \times 10^{-12} \approx 1.9 \times 10^{-11}$, meaning current hardware is approximately 11 orders of magnitude less efficient than the fundamental physical limit.

The discrepancy between measured and theoretical SOPS for each format provides insight into the kernel efficiency beyond simple bandwidth utilization. The efficiency column in the measurement table---defined as the ratio of measured SOPS to the theoretical maximum from the Conservation Law---is near 100\% for high-BPW formats and declines for low-BPW formats. This decline is not a failure of the conservation law but rather a consequence of increased computational overhead per weight element: bit-unpacking, lookup-table indexing, and sign-extension all consume arithmetic cycles that reduce the fraction of time available for memory streaming. The SOPS metric captures this trade-off naturally, unlike FLOPS which would report identical efficiency regardless of format.

\section{SOPS for Training}
\label{sec:sops_training}

The SOPS metric applies to both inference and training phases, though the Info Weight takes on different values depending on the phase. During inference (forward pass only), each weight element participates in one multiply-add (one activation scaled by one weight), yielding $\IW_{\text{forward}} = 32 / \BPW$. During training, each weight element participates in three operations: the forward multiply-add, the backward gradient with respect to the weight, and the backward gradient with respect to the activation. Each of these three operations involves the same weight memory footprint.

For training, the effective Info Weight becomes $\IW_{\text{train}} = 3 \times 32 / \BPW = 96 / \BPW$. This triples the SOPS contribution of each weight element during training relative to inference. However, the time per training step is also larger because the backward pass requires additional computation and memory accesses. The net effect is that training SOPS can be estimated as:

\begin{equation}
\text{SOPS}_{\text{train}} = \frac{E \times \IW_{\text{train}}}{T_{\text{train}} \times 10^{21}}
\label{eq:sops_training}
\end{equation}

where $T_{\text{train}}$ includes both forward and backward passes as well as optimizer updates. The ratio of training SOPS to inference SOPS for the same weight count depends on how $T_{\text{train}}$ scales relative to $\IW_{\text{train}}$.

In practice, training SOPS is typically 3 to 5 times lower than inference SOPS on the same hardware, because the backward pass doubles the arithmetic intensity and the optimizer update adds memory-bound operations. However, because the Info Weight triples during training, the SOPS metric still shows a substantial advantage for quantized formats relative to FP32.

The forward-phase SOPS uses the same formula as inference: $\text{SOPS}_{\text{forward}} = E \times \IW_{\text{forward}} / (T_{\text{forward}} \times 10^{21})$. The backward-phase SOPS is more complex because it involves two distinct computations: computing the gradient with respect to the weight matrix (which reads activations from the forward pass and error signals from the layer above) and computing the gradient with respect to the activations (which reads the weight matrix and error signals). Both operations involve the same weight matrix and therefore the same memory footprint.

For the weight-gradient computation, the memory access pattern mirrors the forward pass: the weight matrix is streamed from memory and multiplied by the activation gradients. The Info Weight for this phase is identical to the forward phase: $\IW_{\text{bw\_weight}} = 32 / \BPW$. For the activation-gradient computation, the weights are read again (or retained in cache if the hardware is large enough), yielding the same Info Weight. Total backward IW is $2 \times 32 / \BPW$, and total training IW is $3 \times 32 / \BPW = 96 / \BPW$.

The optimizer update---typically a weight decay and momentum update---operates on the full-precision gradient accumulators, not on the quantized weights directly. The SOPS contribution of the optimizer update is format-independent and small relative to the forward and backward passes. For large model training, the forward and backward passes dominate the runtime, making the Info Weight framework the primary determinant of training SOPS.

Empirical measurements on the Ryzen 5600GT for training a 64M-parameter model (batch size 8, sequence length 512) yield the following training SOPS values. The OIL4 format achieves approximately 78 pSOPS during training, compared to 353 pSOPS during inference. The training-to-inference ratio of 0.22 reflects both the backward pass overhead and the optimizer update cost. FP32 training achieves approximately 1.9 pSOPS, giving OIL4 a 41x advantage in training throughput---larger than the 8x Info Weight advantage because the smaller memory footprint allows more of the model to fit in cache during training.

The batch size has a significant effect on training SOPS. Larger batches increase the arithmetic intensity (operations per byte loaded) because each weight participates in more forward and backward computations per load. At batch size 1, the arithmetic intensity is minimal and the operation is purely memory-bound, achieving the full Info Weight advantage. At batch size 64, the arithmetic intensity increases by a factor of 64 for the forward pass and 128 for the backward pass, potentially shifting the bottleneck from memory to compute. The SOPS framework captures this naturally: the time $T$ in the denominator grows sublinearly with batch size because the memory access cost is amortized.

The gradient accumulation technique, commonly used to simulate large batch sizes on memory-constrained hardware, interacts with SOPS in an interesting way. Gradient accumulation performs $N$ micro-batches of forward-backward passes, accumulating gradients in FP32 before applying the optimizer update. Each micro-batch incurs the full memory access cost for the weights, so the SOPS per micro-batch is identical to the single-batch case. However, because the optimizer update is performed only once per $N$ micro-batches, the effective SOPS increases by a factor of approximately $N / (N - 1 + \epsilon)$ where $\epsilon$ is the ratio of optimizer update cost to micro-batch cost. For large $N$, this approaches 1, meaning gradient accumulation does not materially affect SOPS.

\section{Comparison with Other Metrics}
\label{sec:comparison_metrics}

Several metrics compete with SOPS for measuring computational throughput. TOPS (Trillions Operations Per Second) is popular in the neural processing unit (NPU) and edge AI communities, where operations are typically integer (INT8) rather than floating-point. GOPS (Giga Operations Per Second) is used for DSP and embedded applications. Each metric has strengths in specific domains, but none captures the information-density dimension that SOPS addresses.

TOPS is the most direct competitor to FLOPS in the AI inference space. An NPU advertising 10 TOPS for INT8 operations may deliver the same raw operation count as a CPU achieving 10 GOPS for FP32 operations, yet the FP32 CPU processes an order of magnitude more information per operation. TOPS inherits FLOPS's blindness to operand precision, substituting integer operation counting for floating-point operation counting without addressing the underlying information-density gap.

The fundamental deficiency shared by FLOPS, TOPS, and GOPS is their lack of a precision-aware normalization factor. All three metrics count operations as atomic, context-independent units. In a system that can switch between FP32, INT8, and binary formats at runtime, a single operation count is meaningless without knowing the precision of each operation. This is analogous to measuring data transfers in packets without specifying the packet size---the metric cannot distinguish between a 64-byte packet and a 1500-byte packet.

SOPS solves this by embedding the precision normalization directly into the metric through the Info Weight. Every operation is weighted by $32 / \BPW$ before being counted, so that one SOPS always represents the same information-theoretic throughput regardless of the underlying format. A system running at 10 SOPS processes the same amount of FP32-equivalent information whether it achieves that throughput through binary weights at high element rates or FP32 weights at low element rates.

\begin{center}
\tablestyle
\begin{tabular}{lcccc}
\toprule
\textbf{Metric} & \textbf{Precision-Aware} & \textbf{Unit Scale} & \textbf{Mixed-Format} & \textbf{Architecture} \\
\midrule
FLOPS & No & $10^9$ & No & General \\
TOPS & No & $10^{12}$ & No & NPU/AI \\
GOPS & No & $10^9$ & No & DSP/Embedded \\
OPS & No & $10^0$ & No & General \\
SOPS & Yes & $10^{21}$ & Yes & Mixed-precision \\
\bottomrule
\end{tabular}
\end{center}

The unit scale of SOPS ($10^{21}$) is substantially larger than FLOPS ($10^9$) or TOPS ($10^{12}$). This is intentional: SOPS is designed to accommodate the full range of mixed-precision throughput from single-core binary inference ($10^{-15}$ SOPS) to hypothetical exascale quantized systems ($10^6$ SOPS and beyond). The sixteen-order-of-magnitude range eliminates the need for multiple coexisting scales and provides a unified framework for comparing fundamentally different systems.

A second advantage of SOPS is its natural support for mixed-format computations. In real-world models, different layers may use different OIL formats (for example, attention layers at OIL4 and feed-forward layers at OIL2). The SOPS contribution of each layer is computed independently and summed: $\text{SOPS}_{\text{total}} = \sum_i E_i \times \IW_i / (T \times 10^{21})$. FLOPS has no mechanism for aggregating heterogeneous precision---one must either report separate FLOPS per layer or compute a meaningless average.

The BOPs (Bit Operations Per Second) metric has been proposed for quantized systems and deserves comparison. BOPs counts the total number of bit-level operations: for a multiply-add of two $b$-bit operands, BOPs counts $b \times b + b$ operations (multiply plus accumulate). While BOPs does capture precision, it conflates operand precision with operation count in a way that does not cleanly separate the information and arithmetic dimensions. SOPS, by contrast, cleanly separates the raw element throughput ($E/T$) from the information weight ($\IW$), allowing each to be analyzed independently.

Another point of comparison is the Roofline model, which visualizes attainable performance as a function of arithmetic intensity. SOPS integrates naturally into the Roofline framework by renormalizing the vertical axis from FLOPS to SOPS. In a SOPS-based Roofline, the memory-bound ceiling is flat regardless of format (because SOPS already incorporates the precision scaling), while the compute-bound ceiling scales with $1/\BPW$. This renormalization makes the Roofline analysis simpler for mixed-precision workloads: one Roofline plot suffices for all formats, rather than requiring separate plots for each precision level. The intersection point---where the memory-bound and compute-bound ceilings cross---shifts leftward (lower arithmetic intensity) for lower-precision formats, reflecting the expanded domain of bandwidth-bound operation.

The TOPS/W (TOPS per Watt) efficiency metric, common in mobile and edge AI, also benefits from the SOPS framework. Replacing TOPS with SOPS yields SOPS/W, a precision-aware efficiency metric that directly measures information throughput per unit energy. For a binary inference engine achieving 10 TOPS/W and an FP32 engine achieving 1 TOPS/W, the comparison in SOPS/W is even more stark: the binary engine's 10 TOPS/W translates to $10 \times 32 = 320$ SOPS/W, dwarfing the FP32 engine's $1 \times 1 = 1$ SOPS/W.

The energy-delay product (EDP), a standard metric for evaluating computational efficiency, also benefits from the SOPS framework. Traditional EDP uses FLOPS or TOPS as the throughput measure, yielding joules-seconds per FLOPS or per TOPS. With SOPS, the EDP becomes joules-seconds per SOPS, capturing both the energy efficiency and the information efficiency of the computation. For OIL4 versus FP32 on the same hardware, the EDP improvement is even larger than the SOPS improvement because the reduced memory traffic also reduces energy consumption---fewer bytes fetched from DRAM means fewer joules expended. Preliminary measurements on the Ryzen 5600GT show that OIL4 achieves a 12x improvement in EDP over FP32, exceeding the 8x SOPS advantage because memory energy dominates the total power budget.

Standardization of SOPS as a benchmark metric faces several challenges. First, SOPS depends on the operation type (GEMM versus element-wise, bandwidth-bound versus compute-bound) in ways that FLOPS does not. A comprehensive SOPS benchmark must specify the kernel type, matrix dimensions, and memory hierarchy configuration to ensure reproducibility. Second, SOPS for a system depends on the specific OIL format mix---a benchmark running at OIL4 cannot be directly compared to one running at OIL2 without normalizing by the Info Weight. Third, the unit scale ($10^{21}$) is unfamiliar to most practitioners and requires sub-SOPS prefixes that are not widely adopted.

\section{Path to 1 SOPS}
\label{sec:path_to_sops}

The goal of one SOPS represents a milestone approximately 11 orders of magnitude beyond current consumer hardware. The memory bandwidth required to sustain 1 SOPS varies quadratically with BPW: OIL4 requires 62.5 EB/s, while BINARY reduces this to 3.9 PB/s---within reach of future HBM4-class multi-chip modules. The quadratic scaling means format density improvements have outsized impact on feasibility.

\begin{center}
\tablestyle
\begin{tabular}{lrr}
\toprule
\textbf{Format} & \textbf{BPW} & \textbf{BW for 1 SOPS} \\
\midrule
BINARY & 1.00 & 3.9\,PB/s \\
TERNARY & 1.58 & 9.8\,PB/s \\
OIL2 & 2.00 & 15.6\,PB/s \\
OIL4 & 4.00 & 62.5\,PB/s \\
OIL8 & 8.00 & 250\,PB/s \\
OIL16 & 16.00 & 1000\,PB/s \\
OIL32 & 32.00 & 4000\,PB/s \\
\bottomrule
\end{tabular}
\end{center}

A concrete scaling roadmap from current hardware to 1 SOPS involves improvements across multiple dimensions. Starting from the Ryzen 5600GT baseline of 0.092 pSOPS at OIL4, each step approximately doubles or quadruples throughput: optimizing FMA efficiency (8.7\% to 80\%) gives 9.2x, AVX512 adds 2x, lower-precision formats add 2x to 8x, and cluster scaling provides the remaining factors.

The transition from OIL4 to TERNARY alone provides a $(4 / 1.58)^2 \approx 6.4$x SOPS gain through quadratic BPW scaling. Cluster scaling provides near-linear gains because each node adds independent memory bandwidth. From the 0.092 pSOPS baseline, optimizing FMA efficiency (9.2x), upgrading to AVX512 (2x), adopting TERNARY (3.2x), and scaling to a 1024-node cluster (128x) yields 0.092 $\times$ 9.2 $\times$ 2 $\times$ 3.2 $\times$ 128 $\approx$ 693 nSOPS. Custom ASICs, optical interconnects, and quantum-assisted sampling bridge the remaining gap from milli-SOPS to 1 SOPS through architectural innovation beyond conventional von Neumann computing.

Despite these challenges, the adoption of SOPS would bring significant benefits to the AI hardware ecosystem. Hardware vendors could advertise SOPS alongside FLOPS for their mixed-precision capabilities, providing a direct comparison across architectures. Cloud providers could bill based on SOPS consumed rather than GPU-hours, aligning cost with the information-processing value delivered. Researchers could compare quantization algorithms in SOPS-normalized terms, separating the format efficiency from the hardware efficiency.

In summary, SOPS does not replace FLOPS, TOPS, or GOPS for their original domains. For homogeneous FP32 computing, FLOPS remains the appropriate metric. For NPU databooks, TOPS is entrenched and sufficient. But for the emerging regime of mixed-precision AI where a single workload spans formats from 1 to 32 bits, SOPS provides the missing informational dimension that makes cross-format comparisons meaningful.
